\section{Results}
\label{sec:results}
Table~\ref{tab:answers} summarises the answers, separating registered findings from post hoc readings; the subsections give the evidence. Subsection~\ref{sec:hosts} and Section~\ref{sec:external} then test how far the answers carry to other hosts, to published instance families and to the standard benchmark.

\begin{table}[t]
\caption{Answers to the three research questions and the check on other hosts (gains in percentage points of the makespan).}
\label{tab:answers}
\centering
{\scriptsize\setlength{\tabcolsep}{3pt}
\begin{tabular}{@{}lp{7.0cm}p{2.8cm}l@{}}
\toprule
 & Answer & Status & Evidence \\
\midrule
RQ1 & Search structure, not the policy (the IG beats the ILS by $0.46\%$ at equal CPU time); budget is $93\%$, $30\%$ or $28\%$ of a local search's gain, by implementation & registered & \ref{sec:rq1} \\
RQ2 & On the synthetic classes no learner beats random choice at equal iterations by more than $0.04$ points; equal-CPU wins are cost effects; a validated action wins & registered & \ref{sec:rq2}, Table~\ref{tab:agents} \\
RQ3 & Headroom is a necessary condition only: outside the discovery set no cell below about $0.05$ points gains beyond the margin, and the standard Q-Learning gains in $29$ of the $79$ above it, but an order-invariant Q-Learning gains in none of the $69$ cells where we repeated it (Section~\ref{sec:order}); a selector trained offline captures the headroom (parity with the validated action, $0.02$--$0.04$; a neural policy, two thirds of the validated gain), no online repair works & registered (held-out test, cross-fitting); the rule and its one-sided reading are exploratory, and the positive half of the rule did not survive the order check & \ref{sec:rq3}, \ref{sec:order}, Fig.~\ref{fig:class}, Table~\ref{tab:selector} \\
Hosts & No online learner gains on a flowshop; on a genetic algorithm only Thompson sampling does, and a validated mixture is never worse; on unrelated machines the validated action beats every learner (also at equal CPU, on published families, on the standard benchmark and at a ten times longer horizon), whose best gets $18\%$ of its gain with uniform times; the apparent $65\%$ with correlated machines was the standard Q-Learning's lock-in on the first action & registered; the mixture was specified after the single action lost & \ref{sec:hosts}, \ref{sec:external}, Table~\ref{tab:hosts} \\
\bottomrule
\end{tabular}}
\end{table}

\subsection{RQ1: the search structure and the budget, not the policy}
\label{sec:rq1}
\paragraph{Structure} Against LPT and seven ILS variants in the twelve classes, the Q-Learning-guided IG is significantly better, after Holm correction, in $72$ of $96$ comparisons (an early Python study, equal iterations; the rest are not significant) and never worse. Yet its advantage over its four ablation siblings (random selection, UCB1 with and without state, Q-Learning without state) is negligible: $0.011\%$ ($90\%$ interval $0.003$ to $0.019$; ten algorithm seeds on $240$ synthetic instances). In the same early study, on two other sets of $4{,}160$ and $2{,}690$ published instances (used nowhere else here), the learned selector beats random selection of the destruction strategy by only $0.009$ to $0.029$ points (all six Holm-adjusted tests significant), and an informative state adds nothing ($-0.006$ points pooled, $p=0.77$). The mechanism is large: with learning removed from both sides, the IG beats the ILS at equal iterations (Python implementation) in $10$ of $12$ classes (the other two differ by less than $0.01\%$) by $0.49\%$, against $0.53\%$ for the learned versions.

This does not depend on how efficiently the ILS is implemented: with vectorised operators ($6.2$ times faster; an iteration costing $0.4$ to $1.9$ times that of the IG), the IG still beats the ILS by $0.46\%$ at equal CPU time ($p_{\mathrm{Holm}}<10^{-80}$), with a positive gain in all twelve classes.

\paragraph{Budget} Adding a short best-improvement descent after each reconstruction improves the IG by $0.150\%$ at an equal number of iterations, but it costs CPU time, so the verdict at equal CPU depends on the cost ratio $K$, and $K$ depends on the implementation. With the first, loop-based descent ($K=9.6$) giving the plain IG the CPU time of the hybrid leaves an advantage of only $0.010\%$: $93\%$ of the gain would have been budget. Vectorising the descent cuts $K$ to $1.8$ and the advantage at equal CPU is $0.105\%$ ($30\%$ budget); compiling both arms gives $K=2.0$ and $0.115\%$ of the $0.159\%$ at equal iterations ($28\%$), a reading registered before the compiled run. For the practical recipe (tuned $d$, start from the best constructive solution) the descent is neutral at equal CPU time ($-0.008\%$ and $-0.006\%$, inside the margin). An evaluation count agrees in sign but not in size ($0.080\%$, $50\%$ budget); equal CPU time remains the budget of choice. The lesson is methodological: the share of a local search's gain that is ``only budget'' ranges from $28\%$ to $93\%$ with the implementation, so a claim about it is a claim about an implementation unless every arm is implemented with comparable efficiency.

\subsection{RQ2: do learners beat random choice, does a static configuration win, and could learning have shown up?}
\label{sec:rq2}
Compiled runs are about a hundred times faster, which allows five seeds per instance in four IG action spaces (Table~\ref{tab:agents}). $K$ in Table~\ref{tab:agents} is the cost of a whole run: $1.01$--$1.08$ for UCB1 and the policy gradient, $1.07$ for Q-Learning with three strategies and $0.50$--$0.60$ in the three larger spaces, because it settles on the cheaper actions (in fact on the first of its list, which is the smallest and cheapest: Section~\ref{sec:order}). \emph{At an equal number of iterations none of the three original learners gains more than $0.016\%$ over random choice in any of the $12$ comparisons}, within the equivalence margin (means over the $12$ classes, Fig.~\ref{fig:class}). \emph{At equal CPU time Q-Learning does win} in the two spaces of sizes, by $0.059\%$ and $0.049\%$, \emph{but only by choosing cheaper actions}: its decision effect is nil ($-0.003$ and $-0.004$ points) and its cost effect is $0.062$ and $0.053$ (a cost effect that an order-invariant version of the learner does not have: Section~\ref{sec:order}). Among the classic AOS methods only the ALNS-style roulette gains: $0.039$ points at equal iterations with five sizes (inside the margin; $0.068$ at equal CPU time, outside it) and nothing at equal iterations with fifteen actions ($0.020$ at equal CPU time); the others, racing aside, stay within $0.013$ points, and online racing, the within-run analogue of a racing configurator, is \emph{worse} than random choice ($-0.012$ and $-0.067$ points). The configuration validated per class is ahead of every one of the nine methods by $0.053$--$0.147$ points in the spaces of $5$ and $15$ actions ($p_{\mathrm{Holm}}<10^{-57}$) and level with the three original learners in the other two, where the oracle headroom is smaller ($0.016$ points behind Q-Learning with three strategies). Where the actions differ in effect and cost (the four ILS operators, $1{,}200$ instances) a DQN and LinUCB lose to the uniform mixture ($0.041\%$ and $0.066\%$), as does the operator validated per class ($0.062\%$).

% Auto-generated by analysis/make_table_cec.py
\begin{table}[t]
\caption{Compiled agents and further online methods in action spaces of $3$ to $30$ actions (P$||C_{\max}$; $12$ classes $\times$ $100$ instances, five seeds, unit = instance mean, $1{,}200$ pairs per row). Gains in percentage points of the makespan with the 90\% interval (positive: the first-named arm is better). Oracle: best fixed action per instance over random choice (one seed per action, hence biased upwards). $K$: CPU time over that of random choice for the same number of iterations (geometric mean of the per-class ratios; equal CPU uses each class's own ratio). Equal iterations: learner over random choice ($T=1000$); equal CPU: learner over random choice with $T=\mathrm{round}(1000K)$; validated: best fixed action per class (tuning instances) over the learner, at equal iterations. The other five methods (probability matching, adaptive pursuit, Thompson sampling, DMAB, racing) are in the text.}
\label{tab:agents}
\centering
{\scriptsize\setlength{\tabcolsep}{3pt}
\resizebox{\textwidth}{!}{\begin{tabular}{@{}lllrlll@{}}
\toprule
Actions & Oracle & Learner & $K$ & Equal iterations & Equal CPU & Validated vs learner \\
\midrule
\multirow{3}{*}{3 strategies} & \multirow{3}{*}{0.046} & Q-Learning & 1.07 & +0.016 [+0.011, +0.021] & +0.011 [+0.006, +0.016] & -0.016 [-0.023, -0.010] \\
 &  & UCB1 & 1.01 & +0.002 [-0.003, +0.007] & -0.000 [-0.005, +0.005] & -0.002 [-0.009, +0.004] \\
 &  & Policy grad. & 1.01 & -0.003 [-0.008, +0.002] & -0.005 [-0.010, -0.000] & +0.003 [-0.003, +0.009] \\ \midrule
\multirow{4}{*}{5 sizes} & \multirow{4}{*}{0.125} & Q-Learning & 0.60 & -0.003 [-0.011, +0.005] & +0.059 [+0.050, +0.068] & +0.095 [+0.087, +0.102] \\
 &  & UCB1 & 1.01 & +0.001 [-0.005, +0.006] & -0.006 [-0.011, -0.001] & +0.091 [+0.085, +0.097] \\
 &  & Policy grad. & 1.03 & -0.000 [-0.005, +0.005] & -0.005 [-0.010, +0.000] & +0.092 [+0.086, +0.098] \\ \cmidrule(l){3-7}
 &  & ALNS roulette & 0.81 & +0.039 [+0.033, +0.045] & +0.068 [+0.062, +0.075] & +0.053 [+0.048, +0.058] \\ \midrule
\multirow{4}{*}{15 (strat.$\times$size)} & \multirow{4}{*}{0.122} & Q-Learning & 0.59 & -0.004 [-0.012, +0.004] & +0.049 [+0.040, +0.058] & +0.084 [+0.076, +0.092] \\
 &  & UCB1 & 1.03 & +0.004 [-0.001, +0.009] & +0.001 [-0.004, +0.005] & +0.076 [+0.070, +0.082] \\
 &  & Policy grad. & 1.06 & +0.000 [-0.005, +0.005] & -0.008 [-0.013, -0.003] & +0.080 [+0.074, +0.086] \\ \cmidrule(l){3-7}
 &  & ALNS roulette & 0.78 & -0.002 [-0.008, +0.003] & +0.020 [+0.014, +0.025] & +0.082 [+0.076, +0.089] \\ \midrule
\multirow{3}{*}{30 (+ local search)} & \multirow{3}{*}{0.049} & Q-Learning & 0.50 & -0.017 [-0.023, -0.011] & +0.016 [+0.010, +0.022] & +0.029 [+0.022, +0.035] \\
 &  & UCB1 & 1.05 & +0.004 [+0.000, +0.007] & +0.002 [-0.001, +0.005] & +0.008 [+0.004, +0.012] \\
 &  & Policy grad. & 1.08 & -0.001 [-0.005, +0.003] & -0.002 [-0.006, +0.001] & +0.013 [+0.008, +0.017] \\
\bottomrule
\end{tabular}}}
\end{table}


\paragraph{Could learning have shown up?} \emph{Power.} With $1{,}200$ pairs the analytic minimum detectable gain at $80\%$ power is $0.006$--$0.014$ points; planting gains in the paired differences gives, in the ten comparisons with a negligible observed effect, $58$--$99.8\%$ power at $0.02$ points and at least $97\%$ at $0.03$, while the largest gain of Q-Learning, UCB1 or the policy gradient is $0.016$: gains of $0.03$ points or more are excluded for them. \emph{Positive controls.} The per-class validated action is detected at $0.092$ and $0.080$ points, and with a reward bonus planted on one action all six learners tested choose it at least $1.5$ times more often than uniform.

\paragraph{Horizon} With five sizes and $T$ from $5{,}000$ to $100{,}000$ the learners stay within $0.006$ points of random choice while the one-seed oracle headroom falls from $0.125\%$ at $T=1{,}000$ to $0.015\%$ at $T=100{,}000$, where the validated action still leads them by $0.012$--$0.017$ points ($p<10^{-12}$). On unrelated machines ($T$ up to $100{,}000$, exploration scaled to $T$ too) the validated gain falls from $1.00$ and $2.18$ ($T=1{,}000$, the same $240$ instances) to $0.31$ and $0.45$ points, Q-Learning never beats random choice ($-0.13$ to $-0.95$) and the validated action beats every learner at every horizon ($0.37$--$2.33$). The longer horizons do not rescue the learner: they shrink the headroom it was supposed to exploit.

\paragraph{Finer times} With $p_j\sim U[1,10^4]$ the cross-fitted headroom is $0.098$ and $0.121$ ($0.142$ and $0.175$ for one seed), yet no learner gains more than $0.015$ points at equal iterations and the validated action beats every learner by $0.105$--$0.136$ points.

\paragraph{A time-varying competitor} The best three-phase schedule of the five sizes per class, chosen on tuning instances, gains $0.095$ against $0.092$ for the validated action on P$||\Cmax$ (parity) and loses $0.166$ to it on unrelated machines ($0.178$ better on the tuning instances: the winner's curse among $125$ candidates): adapting over time adds nothing at the class level.

\paragraph{State-dependent look-ahead} A myopic policy that estimates every action's expected improvement from the actual state at each iteration (eight independent trials per action: up to $120$ extra evaluations per iteration), and a rollout policy built on the validated action (full-horizon rollouts, twelve per action every $50$ iterations), get hundreds of times a learner's evaluations; the validated action beats both in all eight cells ($0.04$--$1.58$ points), and the myopic policy reaches only $6\%$ and $24\%$ of the validated gain on unrelated machines. Knowing exactly the expected one-step reward, which is the signal a Q-Learning agent is trying to learn, does not lead to a better policy than the best fixed action. These policies are references, not oracles or upper bounds: their estimates are noisy and the myopic criterion is not the long-run one.

\paragraph{A configurator} irace~\cite{LopezIbanezEtal2016}, given the tuning cost of the validated action, tunes a continuous destruction size (and strategy) per class: parity with the validated action on P$||\Cmax$ (within $0.012$ points), $0.15$ and $0.19$ points behind on R$||\Cmax$, and still never ahead beyond the margin with ten times the budget. Tuning Q-Learning's four hyperparameters (learning rate, discount, exploration decay and floor) per class selects high exploration floors (median $0.26$--$0.37$) and beats the default Q-Learning on R$||\Cmax$ by $0.58$ and $0.63$ points (paired, on its own test instances), yet the validated action beats it in all four host--space pairs ($0.06$--$1.90$ points). Table~\ref{tab:further} collects these competitors.

% Auto-generated by analysis/make_table_cec5.py
\begin{table}[t]
\caption{Competitors beyond the main battery: gain over random choice at equal iterations, in percentage points of the makespan (positive: better); $1{,}200$ instances per cell ($240$ for the rollout lookahead). The published designs are re-implemented and tuned by us; the lookahead policies are references, not upper bounds. S5, S15: five and fifteen actions. Rows can differ from the paired differences in the text (each is relative to its own reference arm).}
\label{tab:further}
\centering
{\scriptsize\setlength{\tabcolsep}{3pt}
\begin{tabular}{@{}lrrrr@{}}
\toprule
 & \multicolumn{2}{c}{P$||C_{\max}$} & \multicolumn{2}{c}{R$||C_{\max}$} \\ \cmidrule(lr){2-3}\cmidrule(l){4-5}
 & S5 & S15 & S5 & S15 \\
\midrule
Validated single action & +0.092 & +0.080 & +0.947 & +2.078 \\
irace, static configuration & +0.099 & +0.068 & +0.791 & +1.879 \\
irace, tuned Q-Learning & +0.012 & +0.021 & -0.449 & +0.169 \\
Episodic Q-Learning~\cite{KarimiMamaghanEtal2023} & -0.002 & -0.083 & -0.153 & +0.244 \\
FRRMAB-style bandit~\cite{LiFialhoEtal2014} & +0.006 & +0.016 & +0.021 & +0.292 \\
Myopic lookahead (not budget-fair) & +0.019 & +0.035 & +0.058 & +0.497 \\
Rollout lookahead (not budget-fair) & +0.041 & +0.009 & +0.657 & +1.037 \\
\bottomrule
\end{tabular}}
\end{table}


\begin{figure}[t]
\centering
\includegraphics[width=0.78\textwidth]{fig_cec6.pdf}
\caption{Cross-fitted oracle headroom against the gain of the \emph{standard} Q-Learning over random choice at equal iterations (see Section~\ref{sec:order} for the order-invariant learner), per cell: synthetic class by action space ($24$ cells, grey), published family by size tercile by space ($18$ discovery cells, hollow; $18$ held-out, filled), flowshop ($18$ synthetic and $18$ Taillard cells, crosses), genetic algorithm ($24$ cells, triangles) and synthetic class with finer times ($24$ cells, squares). Dotted: the threshold $h^*$. Unrelated-machines cells are off scale.}
\label{fig:class}
\end{figure}

\subsection{RQ3: where does an online selector pay, can that be foretold, and what exploits the headroom?}
\label{sec:rq3}\label{sec:published}\label{sec:selector}
\paragraph{The oracle headroom is mostly noise unless it is cross-fitted} Taken as the best of one seed per action, the headroom of the synthetic classes is $0.125$ and $0.122$ points (five and fifteen actions); choosing the action on two seeds and evaluating it on two fresh ones removes half to two thirds of it ($0.063$ and $0.043$), and $28$--$40\%$ of the headroom of the published instances (only $15$--$20\%$ in family C, middle and upper terciles). On the flowshop, where run-to-run noise is high, it removes $80$--$90\%$ ($0.119\to0.012$ and $0.152\to0.025$). The cross-fitted value is a lower bound, the one-seed value an upper bound (Proposition~\ref{prop:bounds}); on real data the one-seed oracle falls as more runs are used to choose and evaluate on the same runs, and the cross-fitted value lies below all of them (Appendix~\ref{app:estimator}).

\paragraph{Published instances} We ran the nine learners on $2{,}550$ of the published instances, with the hyperparameters and the validated action fixed per family and size tercile on $525$ separate ones. \emph{Here, with the standard implementation, learners did beat random choice} (a reading we withdraw in Section~\ref{sec:order}, where the same held-out instances are rerun with an order-invariant Q-Learning): Q-Learning gains $0.058$ points at equal iterations with fifteen actions and the ALNS roulette $0.057$ with five sizes; in family C, middle and upper size terciles, Q-Learning gains $0.16$ and $0.24$ points, while in family E it \emph{loses} $0.06$--$0.07$. We fixed this reading as hypotheses before running $3{,}950$ held-out instances, retuning nothing. All five confirmatory tests held (Holm over five): Q-Learning gains $0.245$ and $0.149$ points in family C, middle and upper terciles, with fifteen and five actions, and loses $0.052$ and $0.059$ in family E; the ALNS roulette gains $0.056$ with five sizes. At equal CPU time, with its own cost ratio, the validated action is Holm-significantly ahead of the learner in all $54$ learner--space--set pairs (nine learners, two spaces, the synthetic, discovery and held-out sets), though only level (inside the margin) with the ALNS roulette with five sizes ($0.023$ and $0.024$).

\paragraph{A necessary condition, not a sufficient one} Across the cells of a space, the rank correlation between the cross-fitted headroom and Q-Learning's gain is $0.95$ and $0.97$ on the held-out instances (against $0.67$ and $0.60$ with the biased headroom), $0.64$ and $0.66$ across the twelve synthetic classes, and $-0.27$ and $-0.07$ on the flowshop (against $-0.93$ and $-0.78$, a correlation that was mostly noise). We adopted the cross-fitted estimator after the biased one failed in the synthetic classes and the flowshop. These are correlations with the gain of the \emph{standard} Q-Learning; with the order-invariant version, on the $18$ held-out cells, the correlation between the one-seed headroom and the learner's gain is $-0.74$ (Section~\ref{sec:order}): the standard learner gained where the headroom was large because its first action, the cheapest, was a good one for large instances.

\paragraph{One threshold} (This paragraph and the next describe the standard Q-Learning; Section~\ref{sec:order} qualifies what survives of them.) A single threshold, $h^*=0.048$ points, chosen on the $18$ discovery cells and applied unchanged to the $60$ others (held-out, synthetic and flowshop-synthetic cells), classifies $78\%$ of them correctly ($89\%$ held-out, $75\%$ synthetic, $72\%$ flowshop): balanced accuracy $0.80$ (bootstrap over the $30$ classes, each counted in both spaces; $95\%$ interval $0.67$--$0.91$, against $0.60$, $0.42$--$0.79$, with the one-seed headroom; Fig.~\ref{fig:class}). Post hoc checks qualify the rule. Its excess over the $68\%$ of always predicting ``no gain'' is not significant; it does not discriminate within the flowshop (balanced accuracy $0.38$) and classifies only $58\%$ of the cells at finer times (a registered prediction that missed), where Q-Learning loses in the small classes despite a positive headroom; on the genetic algorithm ($24$ cells outside the $60$) it classifies $88\%$.

\paragraph{The one-sided reading} The robust reading is one-sided (post hoc, the $150$ cells outside the discovery set, all hosts): below $h^*$ there are $71$ cells and in none does Q-Learning gain beyond the margin (a point gain above $0.05$ points; $+0.026$ at most); above it there are $79$ others, and it does in $29$, gains something in $42$ and loses in $37$ (a median of $20\%$ of the headroom captured). With margins of $0.02$ and $0.10$ the counts are $2$ and $0$ of $71$ below $h^*$ and $37$ and $15$ of $79$ above it; choosing the action on three seeds instead of two leaves the accuracy unchanged ($0.78$), and so does drawing the random-choice baselines from independent runs (balanced accuracy $0.80$, interval $0.68$--$0.90$). The reading concerns Q-Learning and single-action headroom: a learner that mixes complementary actions can exceed it (Thompson sampling gains $0.051$ in the nine-action genetic algorithm, where the cross-fitted headroom is $-0.073$), and on the flowshop even the validated per-class action exceeds the cross-fitted value. Necessity is nearly a truism for a selector that commits to one action; what is empirical is that the noise-free bound is tight enough to use and far from sufficient.

\paragraph{Horizon in family C} In family C, held-out, middle and upper terciles (every third of the $1{,}020$ instances, three seeds), the cross-fitted headroom falls from $0.222$ points at $T=1{,}000$ to $0.047$ at $T=5{,}000$ and $0.023$ at $T=20{,}000$, Q-Learning's gain from $0.128$ to $0.033$ and $0.017$, and the validated action and Q-Learning become level (differences of $0.002$ and $0.001$ points, $p_{\mathrm{Holm}}>0.9$): the headroom is a short-horizon phenomenon even where it is largest. (Five sizes; on the $340$-instance subsample the gain at $T=1{,}000$ is $0.128$, against $0.149$ on all $1{,}020$ instances above.)

\paragraph{What exploits the headroom: a policy learned offline} We trained an ExtraTrees regressor on $5{,}720$ B, C, E instances disjoint from every tuning and test set, with scale-free features of the instance ($n$, $m$, dispersion and quantiles of the times, bound ratios, the gap of LPT; no family label), to predict the relative performance of each fixed action; at test time it picks one action per instance. The model and the hypotheses were fixed before any test run. On the discovery and held-out instances (Table~\ref{tab:selector}) it gains $0.024$--$0.044$ points over the validated action (inside the margin, hence parity, except in the larger space on held-out instances, where the $90\%$ interval crosses the margin), $0.08$--$0.095$ over Q-Learning and $0.04$--$0.14$ over the ALNS roulette at equal iterations, and captures $0.79$--$0.84$ of the one-seed oracle headroom against $0.58$--$0.66$ for the validated action. A control run afterwards qualifies the result: a per-class lookup table learned on the same instances is level with the selector on the held-out instances, where every class was seen in training ($0.004$ and $0.006$ points, not significant), and the selector wins by only $0.017$ and $0.014$ on the discovery instances ($59\%$ of whose classes were seen): most of the gain is class-level knowledge. On unrelated machines the selector, trained on $1{,}200$ further instances ($18{,}000$ and $54{,}000$ training runs, five times the tuning of the validated action), beats the best online learner of the main battery by $0.72$ and $1.91$ points but captures $60\%$ and $73\%$ of the headroom, no more than the validated action ($-0.056$ and $-0.023$ points; level with the class table). Its training cost is not charged to it above: $85{,}800$ and $257{,}400$ IG runs, about $0.7$ CPU-hours, or $21{,}700$ and $65{,}200$ extra iterations per held-out instance if amortised, a horizon at which the headroom has largely vanished.

% Auto-generated by analysis/make_table_cec2.py
\begin{table}[t]
\caption{Offline instance-conditioned selector (trained on $5{,}720$ disjoint instances, no family information) on the published instances ($2{,}550$ discovery, $3{,}950$ held-out; $T=1000$, five seeds): gain of the selector over each reference, in percentage points of the makespan (positive: the selector is better); at equal CPU time each arm has its own cost ratio $K$; $^\dagger$ not significant after Holm correction. Bottom: share of the oracle headroom captured (oracle of one seed per action, an upper bound; the cross-fitted oracle is a lower bound).}
\label{tab:selector}
\centering
{\scriptsize\setlength{\tabcolsep}{3pt}
\begin{tabular}{@{}lrrrr@{}}
\toprule
 & \multicolumn{2}{c}{S5} & \multicolumn{2}{c}{S15} \\ \cmidrule(lr){2-3}\cmidrule(l){4-5}
 & disc. & held & disc. & held \\
\midrule
vs.\ validated action ($525$ tuning inst.) & +0.024 & +0.026 & +0.027 & +0.044 \\
vs.\ validated action, equal CPU & +0.026 & +0.028 & +0.018 & +0.035 \\
vs.\ static action learned on the same $5{,}720$ & +0.022 & +0.025 & +0.018 & +0.025 \\
vs.\ per-class lookup table, same $5{,}720$ & +0.017 & +0.004$^\dagger$ & +0.014 & +0.006$^\dagger$ \\
vs.\ Q-Learning, equal iterations & +0.078 & +0.090 & +0.081 & +0.095 \\
vs.\ Q-Learning, equal CPU & +0.080 & +0.091 & +0.077 & +0.093 \\
vs.\ ALNS roulette, equal iterations & +0.042 & +0.048 & +0.125 & +0.140 \\
vs.\ ALNS roulette, equal CPU & +0.049 & +0.056 & +0.130 & +0.147 \\
vs.\ random choice & +0.099 & +0.104 & +0.139 & +0.146 \\
\midrule
Headroom captured: validated action & 0.62 & 0.59 & 0.66 & 0.58 \\
Headroom captured: lookup table & 0.69 & 0.75 & 0.74 & 0.80 \\
Headroom captured: selector & 0.83 & 0.79 & 0.82 & 0.84 \\
\bottomrule
\end{tabular}}
\end{table}


\paragraph{Repairing the online learner}\label{sec:repairs} Four attempts do not change the order. Tuning Q-Learning per family ($27$ settings of learning rate, discount and exploration decay) adds $0.009$--$0.014$ points; rewarding plateau moves, as the ALNS score does, adds $0.015$ with five sizes and \emph{costs} $0.007$--$0.009$ with fifteen; the validated action beats both in all eight comparisons, and the selector beats them by $0.06$--$0.10$. A Q table \emph{pre-trained} offline on the same instances, the tabular form of dynamic algorithm configuration, is no better than a cold one with five sizes ($-0.006$ and $+0.005$ points) and much worse with fifteen ($-0.146$ and $-0.143$; after $5.7$ million updates its table is mostly zeros, and breaking ties at random does not help). Optimistic initialisation stays level with random choice ($0.004$--$0.008$ points), is worse than the cold Q-Learning ($0.007$--$0.051$) and loses to the validated action by $0.067$--$0.104$. The limit is the learning signal (bandit feedback), not the cold start: the selector sees every action's outcome on every training instance. (The first three repairs keep the standard tie rule, so they inherit the order sensitivity of Section~\ref{sec:order}.)

\subsection{Other hosts: a flowshop, a genetic algorithm and unrelated machines}
\label{sec:hosts}
We repeated the protocol, in compiled code, on the permutation flowshop (nine classes, $n\in\{20,50,100\}$, $m\in\{5,10,20\}$, with $20$ tuning and $40$ test instances each, plus the $90$ instances ta001--ta090 of Taillard; $T=1000$, three seeds; five and fifteen actions), on a steady-state genetic algorithm for P$||\Cmax$ ($1{,}200$ instances, $T=10{,}000$ offspring, nine and eighteen actions) and on unrelated machines; their gains are in Table~\ref{tab:hosts} and their cells (except the last host's) in Fig.~\ref{fig:class}.

\paragraph{Flowshop} The cross-fitted headroom is only $0.012$ and $0.025$ points. No online learner beats random choice at equal iterations (Q-Learning loses $0.05$--$0.08$ points, and $0.03$--$0.05$ at equal CPU time) and the validated action beats the five learners (Q-Learning, UCB1, the policy gradient, the ALNS roulette, Thompson sampling) in all ten learner--space pairs of the synthetic instances at equal iterations (by $0.020$--$0.104$).

\paragraph{Genetic algorithm} The actions are complementary, so a single validated action can lose to uniform random choice ($-0.057$ points with nine actions). Thompson sampling gains $0.051$ and $0.047$ points, Q-Learning loses $0.096$ and $0.048$, and the competitor that wins is a validated \emph{mixture}, the uniform choice among the three best single actions of each class (specified after the single action lost on this host): it gains $0.050$ and $0.144$ points, beats the validated single action by $0.105$ and $0.067$ and nine of the ten learner--space pairs at equal iterations and at equal CPU time.

\paragraph{Unrelated machines ($R||\Cmax$)} The synthetic classes sit near the ceiling, so we added a host with real headroom, tested with the threshold frozen and the hypotheses registered beforehand (the same twelve classes, $100$ test instances each, $p_{jk}\sim U[1,100]$; predictions fixed after a pilot). The one-seed headroom is $1.5$ and $2.8$ points (cross-fitted $0.43$ and $1.60$), $12$ and $23$ times that of the synthetic P$||\Cmax$ ($7$ and $37$ times cross-fitted), and the validated action gains $0.95$ and $2.08$ points. Yet at equal iterations Q-Learning loses ($-1.06$ and $-0.48$), UCB1 and the policy gradient are at parity with random choice, and, in the main battery, only the optimistic Q-Learning ($+0.17$, five sizes) and the ALNS roulette ($+0.15$, fifteen actions) gain beyond the margin, at most $18\%$ and $7\%$ of the validated action's gain; the validated action beats every learner in all ten learner--space pairs, by $0.77$--$2.52$ points.

\emph{At equal CPU time} Q-Learning with fifteen actions (and, not significantly, with five: $0.10$) and the ALNS roulette do gain over random choice ($0.58$ and $0.39$--$0.45$ points), but through the cost effect (decision effects $-0.48$ and $-0.09$ to $+0.15$), and the validated action, with its own cost ratio, still beats all five learners in all ten pairs ($0.94$--$2.17$).

\emph{Two published AOS designs}, tuned per space on the tuning instances as their hyperparameters were not available (an episodic Q-Learning with a binary state in the style of~\cite{KarimiMamaghanEtal2023} and an FRRMAB-style sliding-window bandit~\cite{LiFialhoEtal2014}), do not close the distance (Table~\ref{tab:further}): neither gains beyond the margin on P$||\Cmax$, and on R$||\Cmax$ they reach at most $14\%$ of the validated action's gain, which beats both in all eight comparisons ($0.06$--$1.82$ points). The frozen rule is not a classifier here ($42\%$ of the $24$ cells: $19$ lie above $h^*$ and Q-Learning gains in only $5$ of those), but its one-sided reading holds out of sample: none of the $5$ cells below $h^*$ shows a gain beyond the margin.

% Auto-generated by analysis/make_table_cec2.py
\begin{table}[t]
\caption{Three further hosts: gain over random choice in percentage points of the makespan (positive: better). Permutation flowshop (F5, F15: five and fifteen actions; $360$ synthetic and $90$ Taillard instances). Unrelated machines (R$||C_{\max}$; $T=1000$; $1{,}200$ instances; R5, R15: five and fifteen actions; the threshold and the hypotheses were frozen before the test, and the optimistic Q-Learning is among the other learners, whereas the published designs and the configurator are reported in the text). Genetic algorithm (G9, G18: nine and eighteen actions). The cross-fitted oracle chooses the action on two seeds and evaluates it on two others (four seeds per fixed action); `--' where not run. Rounded rows can differ from the paired differences in the text (each is relative to its own reference arm).}
\label{tab:hosts}
\centering
{\scriptsize\setlength{\tabcolsep}{4pt}
\resizebox{\textwidth}{!}{\begin{tabular}{@{}lrrrrrrrr@{}}
\toprule
 & \multicolumn{2}{c}{Flowshop, synth.} & \multicolumn{2}{c}{Flowshop, Taillard} & \multicolumn{2}{c}{Unrelated machines} & \multicolumn{2}{c}{Genetic alg.} \\ \cmidrule(lr){2-3}\cmidrule(lr){4-5}\cmidrule(lr){6-7}\cmidrule(l){8-9}
 & F5 & F15 & F5 & F15 & R5 & R15 & G9 & G18 \\
\midrule
Oracle headroom, one seed & +0.119 & +0.152 & +0.132 & +0.186 & +1.498 & +2.806 & +0.097 & +0.188 \\
Oracle headroom, cross-fitted & +0.012 & +0.025 & -0.006 & +0.024 & +0.428 & +1.603 & -0.073 & +0.070 \\
Validated single action & +0.028 & +0.043 & +0.012 & +0.031 & +0.947 & +2.078 & -0.057 & +0.076 \\
Validated mixture (top 3) & -- & -- & -- & -- & -- & -- & +0.050 & +0.144 \\
Q-Learning, equal iterations & -0.050 & -0.061 & -0.059 & -0.080 & -1.059 & -0.484 & -0.096 & -0.048 \\
Best other learner (main battery), equal iter. & +0.008 & +0.009 & +0.006 & +0.005 & +0.174 & +0.148 & +0.051 & +0.047 \\
Q-Learning, equal CPU & -0.036 & -0.047 & -0.030 & -0.037 & +0.100 & +0.584 & -0.113 & -0.054 \\
Offline selector & +0.036 & +0.050 & +0.034 & +0.047 & +0.895 & +2.059 & -- & -- \\
\bottomrule
\end{tabular}}}
\end{table}


\subsection{A caveat on the standard Q-Learning: it depends on the order of the actions}
\label{sec:order}
The standard Q-Learning of our batteries behaved in ways that looked like learning but were not. On the standard benchmark it lost to random choice by up to $1.19$ points in cells where the static choice gains most (Section~\ref{sec:fanjul}); on the published families of unrelated machines it gained up to $65\%$ of the validated gain (Section~\ref{sec:families}); and at equal CPU time it won ``by choosing cheaper actions'' (Section~\ref{sec:rq2}). A diagnostic of the standard benchmark (Entry~258, registered before it was run) showed what the learner actually does: it puts $55$--$73\%$ of its choices on the action listed first (random choice: $1/A$, between $3\%$ and $20\%$), while UCB1 and the policy gradient stay at the share of random choice ($3$--$20\%$). Table~\ref{tab:order_a} gives the shares.

% Auto-generated by analysis/make_table_journal.py
\begin{table}[t]
\caption{Order sensitivity of the standard Q-Learning on the standard benchmark ($140$ test instances, one per cell of type, $m$ and $n$; two seeds; gains over random choice in points, positive = better). ql: standard; ql\_v0: the same learner with the class's validated action moved to index $0$; val: the validated action. Index-0 share: fraction of the learner's choices that fall on the action listed first (random choice: $1/A$); ql\_v0 share: fraction on the validated action.}
\label{tab:order_a}
\centering
{\scriptsize\setlength{\tabcolsep}{5pt}
\begin{tabular}{@{}lrrrrrr@{}}
\toprule
Space & $1/A$ & ql index-0 share & ql & ql\_v0 & val & ql\_v0 share of val \\
\midrule
B5 & 0.20 & 0.73 & $-$0.36 & 0.73 & 0.90 & 0.74 \\
B15 & 0.07 & 0.71 & $-$0.02 & 1.30 & 1.62 & 0.70 \\
B30 & 0.03 & 0.55 & $-$0.54 & 2.02 & 2.35 & 0.68 \\
\bottomrule
\end{tabular}}
\end{table}


The first causal test, reversing the order of the actions, was badly designed: it made the learner lock on the last action of the list, another extreme and poor action, so it did not help (three of the seven predictions of that stage held). The second test moves the validated action of each class to index $0$ without changing anything else. The standard Q-Learning then gains $0.73$, $1.30$ and $2.02$ points over random choice in the three spaces, $81\%$, $81\%$ and $86\%$ of the validated gain (all five registered predictions of this stage held), and it plays the validated action in $68$--$74\%$ of its choices. Standing on the first index, it recovers most of the best fixed action; standing on any other, it recovers an average one. The learner's results depend on where the good action was listed.

\paragraph{The order-invariant learner.} We repeated the cells where the standard learner seemed to gain or lose with a version that permutes the actions at random in each run (\textsc{ql-perm}; Table~\ref{tab:order_b}). On the standard benchmark it gains over random choice beyond the margin in none of the $21$ cells (mean $-0.30$ points, against $-0.25$ for the standard learner); on the three published families it loses to random choice in all six cells (from $-0.05$ to $-0.66$ points; the standard learner's gains of $0.67$ and $0.83$ points with correlated machines were $53\%$ and $65\%$ of the validated gain, those of the permuted learner are negative), and the validated action beats it in all six; on the synthetic $P||C_{\max}$ battery it gains in none of the $24$ cells and the validated action beats it in $23$. Four of the six registered predictions on the standard benchmark and the families held, and all four on $P||C_{\max}$.

% Auto-generated by analysis/make_table_journal.py
\begin{table}[t]
\caption{The order-invariant Q-Learning (ql\_perm: a random permutation of the actions in each run) against the standard one and the validated action: mean gain over random choice across the cells of each battery (points, positive = better) and number of cells with a gain beyond the equivalence margin. P$||C_{\max}$: $12$ classes $\times$ spaces S5 and S15.}
\label{tab:order_b}
\centering
{\scriptsize\setlength{\tabcolsep}{4pt}
\begin{tabular}{@{}lrrrrr@{}}
\toprule
Battery & cells & ql & ql\_perm & val & cells with ql\_perm gain \\
\midrule
Standard benchmark (Fanjul-Peyro, Ruiz) & 21 & $-$0.25 & $-$0.30 & 1.60 & 0 \\
Published $R||C_{\max}$ families & 6 & 0.28 & $-$0.33 & 1.28 & 0 \\
Synthetic $P||C_{\max}$ & 24 & 0.00 & $-$0.06 & 0.09 & 0 \\
\bottomrule
\end{tabular}}
\end{table}


\paragraph{The cost effect.} In the $P||C_{\max}$ battery the actions are listed by destruction size, so the first action is the cheapest. The standard learner evaluates $0.44$--$0.46$ times as many solutions as random choice per run; the order-invariant learner evaluates $0.99$--$1.01$ times as many. The statement that at equal CPU time Q-Learning wins only by choosing cheaper actions is therefore true of the standard implementation, and its cause is the lock-in on the first index, not a policy learned from the reward.

\paragraph{What this changes.} It does not change the main conclusion; it makes it stronger: an online Q-Learning does not learn a useful policy in these budgets, and the largest gains attributed to it came from a good action standing first in the list. It changes three things in how the result must be read. First, ``Q-Learning gains in four of six cells of the published families'' is an artefact of the order of the actions and should not be read as evidence that online learning pays there. Second, the protocol needs a rule it did not have: report the learner's order sensitivity (here by permuting the actions) or break ties at random. Third, in the earlier batteries the standard learner performs like a single fixed action placed arbitrarily in the list; the optimistic Q-Learning, UCB1 and the other bandits, which do not repeat the first index, are the more reliable representatives of online learning and are the ones the claims rest on.


\section{Published instance families and the standard benchmark}
\label{sec:external}
The synthetic classes and our own generators are the weakest point of any such study, because they were chosen by the authors. We therefore repeated the protocol on instances that others defined, in two steps of increasing external validity: three published families of $R||\Cmax$ instances, for which we kept the frozen learners and re-chose only the validated action (Section~\ref{sec:families}); and the standard benchmark of \citet{FanjulPeyroRuiz2010}, which comes with the best known solutions of three exact methods run by third parties and lets us report the distance to them (Section~\ref{sec:fanjul}).

\subsection{Three published families of $R||\Cmax$}
\label{sec:families}
\paragraph{Design.} The families are the uncorrelated one ($p_{jk}\sim U\{10,100\}$), the one with correlated jobs and the one with correlated machines (Section~\ref{sec:instances}), with $25$ sizes each ($m\in\{3,5,10,15,20\}$, $n\in\{30,50,80,100,200\}$), $20$ tuning and $20$ test instances per class, hence $500$ test instances per family and space. Everything that depends on the generator is re-estimated on the tuning instances (the validated action per class); the learners keep the hyperparameters frozen from the synthetic batteries, so the comparison cannot favour them by tuning to the new generator, nor penalise them by the lack of it. The unit of analysis is the mean over seeds of each instance, the cells are (family, space), and the predictions $X_1$--$X_6$ were recorded before the test run.

\paragraph{Results.} The cross-fitted headroom exceeds the frozen threshold in all six cells ($0.52$ to $2.31$ points; Table~\ref{tab:families}), so the necessary condition of the rule of use is met everywhere; the question is whether it is sufficient. The answer is mixed. The validated action beats every learner in all $30$ learner--cell pairs ($0.42$--$2.64$ points) and no learner beats it beyond the margin, as in every previous battery. But Q-Learning gains over random choice beyond the margin in four of the six cells (uncorrelated, correlated jobs and correlated machines with fifteen actions, and correlated machines with five), up to $0.67$ and $0.83$ points with correlated machines, that is $53\%$ and $65\%$ of the validated gain; with five actions it loses $0.45$ points in the uncorrelated family and is level with random choice in correlated jobs ($-0.01$). Of the six registered predictions of this battery four held. The two that did not were that Q-Learning would gain in no cell (it gained in four) and that the best learner would obtain at most $30\%$ of the validated gain wherever the validated action gains (it obtained up to $65\%$). We first read this as the most favourable evidence for online learning that we found. We withdraw that reading in Section~\ref{sec:order}: the gains are an artefact of the order of the actions (the standard learner repeats the first action of its list), and an order-invariant Q-Learning loses to random choice in all six cells.

\paragraph{Stronger competitors.} We then gave the learners the same tuning effort as the validated action and added the two published designs of Section~\ref{sec:learners}. irace received, as its budget, the cost of the tuning of the validated action ($200$ and $600$ experiments per class); it was run to find (a) a static configuration (a continuous destruction fraction and, for fifteen actions, the strategy) and (b) the four hyperparameters of Q-Learning. The static configuration of irace is at parity with the validated action in every cell (differences of at most $0.07$ points; no cell where it beats the validated action beyond the margin, as predicted in $Y_1$). The irace-tuned Q-Learning does better than the default one and gains beyond the margin in five of the six cells (predicted: at most two), up to $0.82$ points (it keeps the tie rule of the standard learner, so the caveat of Section~\ref{sec:order} applies to it as well); the two published designs gain up to $0.54$ points. The validated action beats the three in all $18$ comparisons ($0.40$--$2.19$ points), as predicted ($Y_3$, $Y_4$). Of the six predictions of this battery five held. Thus, even with the tuning of irace and with designs from the literature, a static choice validated on $20$ tuning instances per class is at least level with every alternative in every cell.

% Auto-generated by analysis/make_table_journal.py
\begin{table}[t]
\caption{Gains over random choice (points of makespan, positive = better) on the three published families of $R||C_{\max}$ (generated from the description of \citet{AbladEtAl2021}; $500$ test instances per family and space, $T=1000$ iterations). $h_{\mathrm{cf}}$: cross-fitted headroom; val: validated static action; cfg: irace-configured static action; ql, ucb, pg, alns, qo: Q-Learning, UCB1, policy gradient, ALNS roulette, optimistic Q-Learning; q: irace-tuned Q-Learning; iq, fr: the two published AOS designs. Bold: gain over random choice confirmed (Holm within method) and not equivalent to zero at the margin $\pm0.05$. No column beats the validated one beyond the margin; irace's static action is at parity with it.}
\label{tab:families}
\centering
{\scriptsize\setlength{\tabcolsep}{3.5pt}
\resizebox{\textwidth}{!}{%
\begin{tabular}{@{}llrrrrrrrrrrr@{}}
\toprule
Family & Actions & $h_{\mathrm{cf}}$ & val & cfg & ql & q & ucb & pg & alns & qo & iq & fr \\
\midrule
uncorrelated & 5 & 0.88 & 1.15 & \textbf{1.09} & $-$0.45 & $-$0.05 & $-$0.02 & 0.01 & \textbf{0.11} & $-$0.11 & $-$0.12 & \textbf{0.05} \\
correlated jobs & 5 & 0.52 & 0.55 & \textbf{0.54} & $-$0.01 & \textbf{0.14} & $-$0.01 & 0.00 & \textbf{0.12} & $-$0.06 & \textbf{0.04} & \textbf{0.05} \\
correlated machines & 5 & 1.18 & 1.26 & \textbf{1.28} & \textbf{0.67} & \textbf{0.70} & 0.00 & 0.01 & \textbf{0.50} & \textbf{0.15} & $-$0.03 & \textbf{0.30} \\
uncorrelated & 15 & 2.31 & 2.59 & \textbf{2.59} & \textbf{0.46} & \textbf{0.75} & 0.09 & 0.02 & $-$0.06 & \textbf{0.42} & \textbf{0.54} & \textbf{0.41} \\
correlated jobs & 15 & 0.79 & 0.85 & \textbf{0.82} & \textbf{0.20} & \textbf{0.35} & 0.02 & $-$0.01 & $-$0.08 & \textbf{0.04} & \textbf{0.13} & \textbf{0.07} \\
correlated machines & 15 & 1.37 & 1.29 & \textbf{1.24} & \textbf{0.83} & \textbf{0.82} & \textbf{0.08} & 0.04 & \textbf{0.28} & \textbf{0.20} & $-$0.31 & \textbf{0.25} \\
\bottomrule
\end{tabular}}}
\end{table}


\subsection{The standard benchmark of Fanjul-Peyro and Ruiz}
\label{sec:fanjul}
The benchmark of \citet{FanjulPeyroRuiz2010} (Section~\ref{sec:instances}) has $1{,}400$ instances of seven types (uniform times in $[1,100]$, $[10,100]$, $[100,120]$, $[100,200]$ and $[1000,1100]$, correlated jobs and correlated machines) with $m\in\{10,20,30,40,50\}$ and $n\in\{100,200,500,1000\}$, and best known values from CPLEX, Gurobi and an exact Lagrangian method run by third parties. We used the deterministic split of the repository's manifest (per cell of type, $m$ and $n$: two instances for tuning, five for testing, three reserved): $280$ tuning and $700$ test instances, $100$ test instances per type and action space. This is the host with the largest instances of the study, up to $n=1000$ jobs on $50$ machines.

\paragraph{Design.} The protocol is the compiled one of the earlier batteries, and nothing is tuned to this benchmark except the validated action. A class is a pair (type, $n$), $28$ in all; the validated action of a class is the best fixed action on the $10$ tuning instances of the class ($2$ seeds each). Three action spaces cover the same range as before, with the destruction sizes expressed as fractions of $n$ ($1$, $2$, $5$, $10$ and $20\%$, at least two jobs) so that the grid has the same meaning at $n=100$ and $n=1000$: B5 (five sizes of random destruction), B15 (three destruction strategies $\times$ five sizes) and B30 (B15 without and with a local search after each reconstruction). The seven learners are those of Section~\ref{sec:learners} (Q-Learning, UCB1, the policy gradient, the ALNS roulette, the optimistic Q-Learning and the two published AOS designs), with the hyperparameters frozen from the earlier batteries (those of the five- and fifteen-action synthetic spaces; B30 uses the fifteen-action values). Each method runs $T=1000$ iterations with three seeds per instance, and the unit of analysis is the mean over seeds. The cells are (type, space), $21$ in all, and the seven predictions $Z_1$--$Z_7$ and the hashes of the code and of the validated actions were recorded before the test run.

\paragraph{The validated action.} It gains over random choice beyond the margin in all $21$ cells ($Z_7$; the gain ranges from $0.14$ to $5.08$ points, Table~\ref{tab:fanjul}), and it is better than every learner in all $147$ learner--cell pairs ($Z_1$: $21$ of $21$ against Q-Learning), with no learner beating it beyond the margin in any of them ($Z_2$). In B30 the validated action is a local-search action in all $28$ classes. Measured against the best known solutions of the $700$ test instances, its mean gap is $8.32\%$ in B5, $8.32\%$ in B15 and $4.89\%$ in B30, against $9.35\%$, $10.21\%$ and $7.42\%$ for random choice; the local search helps absolute quality ($Z_4$ and $Z_5$ held, $Z_4$ narrowly at $4.89\%$ against the registered $5\%$). By type the gap of the validated action in B30 runs from $0.5\%$ ($U[1000,1100]$) to $9.9\%$ ($U[10,100]$). These gaps are those of a plain IG run for a thousand iterations, not of a method designed for this benchmark, and the protocol compares arms, not absolute quality.

\paragraph{The learners.} The learners gain over random choice beyond the margin in $39$ of the $147$ learner--cell pairs, and unevenly: Q-Learning in $3$ of its $21$ cells (all in B15, with gains of $0.07$, $0.15$ and $0.39$ points), UCB1 in one, the policy gradient in none, the ALNS roulette in six (all in B30), the optimistic Q-Learning in ten, the episodic design in seven and the FRRMAB-style design in twelve. The largest fraction of the validated gain recovered by a learner is $0.68$ (the FRRMAB-style design on $U[1000,1100]$ with five actions, where the validated gain is only $0.14$ points); in every other cell it is at most $0.46$. Two of the seven registered predictions failed. $Z_3$, that Q-Learning gains beyond the margin in at least one type in each space, failed because it gains in no type in B5 or B30. $Z_6$, that the cells where the validated action gains most are those where Q-Learning gains most (a Spearman correlation of at least $0.5$ over the $21$ cells), failed with the opposite sign: $-0.46$ ($-0.75$ in B5, $+0.14$ in B15 and $-0.68$ in B30, post hoc).

\paragraph{Reading.} The headroom here is not in doubt: the validated action gains at least $0.14$ points in every cell, far above the threshold of $0.05$, and by up to $5.08$ points. Yet the standard Q-Learning gains beyond the margin in three cells of twenty-one, loses to random choice by up to $1.19$ points, and does worst where the static choice matters most (B30, $-0.94$ points on $U[100,200]$ and $-0.88$ on correlated machines). The headroom is a necessary condition and, as in the earlier hosts, not a sufficient one; in this benchmark the relation that held on the published families of Section~\ref{sec:families} (the learner gains most where the headroom is largest) is not reproduced. We did not investigate the mechanism. One post hoc observation: in B30 the validated action evaluates $2.03$ times as many solutions per run as random choice, and Q-Learning $0.72$ times, so it concentrates on the cheaper actions while the validated action is a local-search action. The registered predictions of this battery held in five of seven cases.

% Auto-generated by analysis/make_table_journal.py
\begin{table}[t]
\caption{Standard benchmark of Fanjul-Peyro and Ruiz: gain over random choice (points of makespan, positive = better; $100$ test instances per type and space, $T=1000$ iterations, three seeds) of the validated action (val), of Q-Learning (ql) and of the best of the seven learners in each cell (best; the learner is named), and the mean gap of the validated action to the best known solution (\%). Bold: gain over random choice confirmed (Holm within learner, $21$ cells) and not equivalent to zero at the margin $\pm0.05$. B5: five destruction sizes; B15: three strategies $\times$ five sizes; B30: B15 without and with local search.}
\label{tab:fanjul}
\centering
{\scriptsize\setlength{\tabcolsep}{3pt}
\resizebox{\textwidth}{!}{%
\begin{tabular}{@{}lrrrlrrrrlrrrrlr@{}}
\toprule
 & \multicolumn{5}{c}{B5} & \multicolumn{5}{c}{B15} & \multicolumn{5}{c}{B30} \\
\cmidrule(lr){2-6}\cmidrule(lr){7-11}\cmidrule(lr){12-16}
Type & val & ql & best & & gap & val & ql & best & & gap & val & ql & best & & gap \\
\midrule
$U[1,100]$ & 2.10 & $-$1.19 & 0.42 & qo & 10.4 & 4.22 & $-$0.37 & \textbf{0.54} & fr & 10.3 & 5.08 & $-$0.57 & \textbf{0.78} & qo & 6.8 \\
$U[10,100]$ & 1.61 & $-$0.41 & 0.18 & pg & 14.9 & 3.27 & \textbf{0.39} & \textbf{0.66} & iq & 14.8 & 4.15 & $-$0.42 & \textbf{0.61} & qo & 9.9 \\
$U[100,120]$ & 0.27 & $-$0.15 & \textbf{0.12} & fr & 3.6 & 0.60 & 0.06 & \textbf{0.25} & iq & 3.6 & 0.51 & $-$0.17 & 0.10 & qo & 2.1 \\
$U[100,200]$ & 0.84 & $-$0.27 & \textbf{0.38} & qo & 11.3 & 1.56 & \textbf{0.15} & \textbf{0.32} & iq & 11.2 & 2.93 & $-$0.94 & \textbf{0.39} & alns & 4.8 \\
$U[1000,1100]$ & 0.14 & $-$0.03 & \textbf{0.10} & fr & 1.7 & 0.43 & \textbf{0.07} & \textbf{0.14} & fr & 1.7 & 0.40 & $-$0.06 & \textbf{0.09} & qo & 0.5 \\
correlated jobs & 0.24 & $-$0.16 & 0.02 & fr & 4.3 & 0.44 & $-$0.04 & 0.02 & fr & 4.3 & 0.50 & $-$0.20 & \textbf{0.13} & alns & 3.3 \\
correlated machines & 1.09 & $-$0.07 & \textbf{0.29} & fr & 12.0 & 0.94 & 0.08 & \textbf{0.26} & fr & 12.3 & 2.22 & $-$0.88 & \textbf{0.63} & alns & 6.9 \\
\bottomrule
\end{tabular}}}
\end{table}



